\documentclass{article}
\usepackage{amsmath,amssymb}
\usepackage{geometry}
\geometry{margin=1in}

\begin{document}

\title{CSE400 – Fundamentals of Probability in Computing\\
Lecture L17: Joint Random Variables and Joint Distributions}
\date{}
\maketitle

\section{Introduction: Why Study Joint Random Variables?}

In many real-world situations, we are interested in two or more random variables simultaneously rather than analyzing a single random variable. Studying such variables allows us to understand relationships between different stochastic quantities.

\subsection*{Real-World Examples}

Examples where two random variables must be studied together include:

\begin{itemize}
\item Height and weight of giraffes
\item IQ and birthweight of children
\item Frequency of exercise and rate of heart disease
\item Air pollution level and respiratory illness rate
\item Number of Facebook friends and age of members
\end{itemize}

The purpose of studying these variables jointly is to determine how they are related.

\subsection{Example Applications}

\subsubsection*{Smart Transportation System}

Two random variables are defined:

\begin{itemize}
\item $X$: Vehicle Speed
\item $Y$: Traffic Density
\end{itemize}

Key question:

Can vehicle speed be high when traffic density is also high?

The answer may vary depending on conditions, indicating that $X$ and $Y$ are correlated random variables.

Applications include:

\begin{itemize}
\item Adaptive traffic signals
\item Congestion prediction
\item Autonomous vehicle routing
\end{itemize}

\subsubsection*{Wireless Network Performance}

Two random variables are considered:

\begin{itemize}
\item $X$: Signal-to-Noise Ratio (SNR)
\item $Y$: Data Throughput
\end{itemize}

Observation:

Throughput depends on SNR but may also depend on congestion.

Question:

If SNR is high, will throughput always be high?

Answer: Not necessarily.

Applications:

\begin{itemize}
\item 5G scheduling
\item Adaptive modulation
\item Network planning
\end{itemize}

\subsubsection*{Additional Examples}

Other scenarios involving joint random variables include:

\begin{itemize}
\item Weather prediction

$X$: Rainfall

$Y$: Temperature

\item Drone delivery

$X$: Wind speed

$Y$: Delivery time

\item Rolling two dice

$X$: value of one die

$Y$: value of another die
\end{itemize}

\subsection{Purpose of Joint Distributions}

When random variables are considered together:

\begin{itemize}
\item They have a joint distribution.
\item This distribution allows computation of probabilities involving both variables.
\item It helps analyze the relationship between the variables.
\end{itemize}

Special cases:

\begin{itemize}
\item If variables are independent, analysis is simpler.
\item If variables are dependent, measures such as covariance and correlation are used.
\end{itemize}

\section{Discrete Case: Joint Probability Mass Function (Joint PMF)}

\subsection{Setup}

Suppose two discrete random variables:

\[
X = \{x_1,x_2,\ldots,x_n\}
\]

\[
Y = \{y_1,y_2,\ldots,y_m\}
\]

The ordered pair $(X,Y)$ takes values from the product set

\[
\{(x_1,y_1),(x_1,y_2),\ldots,(x_n,y_m)\}
\]

\subsection{Definition: Joint PMF}

The joint probability mass function is defined as

\[
p(x_i,y_j)=P(X=x_i,Y=y_j)
\]

Interpretation:

This represents the probability that both events occur simultaneously, meaning:

\begin{itemize}
\item $X$ takes value $x_i$
\item $Y$ takes value $y_j$
\end{itemize}

\subsection{Joint PMF Table Representation}

The joint PMF is commonly represented in a table format.

Rows correspond to values of $X$ and columns correspond to values of $Y$.

Example table structure:

\[
\begin{array}{c|ccccc}
X \backslash Y & y_1 & y_2 & \dots & y_j & \dots y_m \\
\hline
x_1 & p(x_1,y_1) & p(x_1,y_2) & \dots & p(x_1,y_j) & \dots p(x_1,y_m) \\
x_2 & p(x_2,y_1) & p(x_2,y_2) & \dots & p(x_2,y_j) & \dots p(x_2,y_m) \\
\vdots & \vdots & \vdots & \vdots & \vdots & \vdots
\end{array}
\]

\subsubsection*{Key Properties}

Probabilities are bounded:

\[
0 \le p(x_i,y_j) \le 1
\]

Total probability equals one:

\[
\sum_{i=1}^{n}\sum_{j=1}^{m} p(x_i,y_j)=1
\]

\section{Example: Rolling Two Dice}

\subsection*{Experiment}

Roll two fair dice.

Define random variables:

\begin{itemize}
\item $X$: value on the first die
\item $T$: total sum of both dice
\end{itemize}

Each ordered outcome $(i,j)$ occurs with probability:

\[
\frac{1}{36}
\]

Therefore, we construct the joint probability table for $(X,T)$.

\subsection*{Observations from the Joint Table}

From the table shown in the lecture slides:

\begin{itemize}
\item Entries are either $0$ or $\frac{1}{36}$.
\item Variables $X$ and $T$ are not independent.
\end{itemize}

\section{Continuous Case: Joint Probability Density Function (Joint PDF)}

The continuous case is analogous to the discrete case.

\subsection*{Correspondence Between Cases}

\begin{tabular}{ll}
Discrete & Continuous \\
Discrete values & Continuous intervals \\
Joint PMF & Joint PDF \\
Sums & Double integrals
\end{tabular}

If:

\[
X\in[a,b],\quad Y\in[c,d]
\]

then

\[
(X,Y)\in[a,b]\times[c,d]
\]

\subsection{Definition: Joint PDF}

A function $f(x,y)$ is called the joint probability density function of $X$ and $Y$ if

\[
P((X,Y)\text{ lies in small rectangle}) \approx f(x,y)\,dx\,dy
\]

Thus, the probability over a small region equals:

\[
f(x,y)\,dx\,dy
\]

\subsection{Key Properties}

Non-negativity

\[
f(x,y)\ge0
\]

Normalization condition

\[
\int_c^d\int_a^b f(x,y)\,dx\,dy = 1
\]

Important note:

A density value $f(x,y)$ may be greater than $1$, since it represents a density, not a probability.

\section{Worked Example}

Given joint PDF:

\[
f(x,y)=4xy
\]

for

\[
0\le x\le1,\quad 0\le y\le1
\]

The pair $(X,Y)$ lies in the unit square

\[
[0,1]^2
\]

\subsection*{Task 1: Verify Valid Joint PDF}

Compute the normalization integral.

\[
\int_0^1\int_0^1 4xy\,dx\,dy
\]

Step 1:

\[
\int_0^1 \left(\int_0^1 4xy\,dx\right)dy
\]

Step 2:

\[
\int_0^1 2y\,dy
\]

Step 3:

\[
[y^2]_0^1 = 1
\]

Thus:

\[
f(x,y)\ge0
\]

total probability = 1

Therefore, $f(x,y)$ is a valid joint PDF.

\subsection*{Task 2: Compute Probability of Event}

Event:

\[
A=\{X<0.5,Y>0.5\}
\]

Compute:

\[
P(A)=\int_0^{0.5}\int_{0.5}^{1}4xy\,dy\,dx
\]

Step 1:

\[
\int_0^{0.5}[2xy^2]_{0.5}^{1}dx
\]

Step 2:

\[
\int_0^{0.5}2x(1-0.25)dx
\]

Step 3:

\[
\int_0^{0.5}\frac{3}{2}x\,dx
\]

Final result:

\[
P(A)=\frac{3}{16}
\]

\section{Joint Cumulative Distribution Function (Joint CDF)}

Suppose $X$ and $Y$ are jointly distributed random variables.

Notation:

\[
X\le x,\quad Y\le y
\]

represents the event:

\[
\{X\le x \text{ and } Y\le y\}
\]

\subsection*{Definition}

The joint cumulative distribution function is defined as

\[
F(x,y)=P(X\le x,Y\le y)
\]

\section{Joint CDF for Continuous Variables}

If $X$ and $Y$ are continuous random variables with joint PDF $f(x,y)$, then

\[
F(x,y)=\int_c^y\int_a^x f(u,v)\,du\,dv
\]

\subsection*{Recovering the Joint PDF}

The joint PDF can be obtained using partial derivatives:

\[
f(x,y)=\frac{\partial^2}{\partial x\partial y}F(x,y)
\]

\section{Joint CDF for Discrete Variables}

If $X$ and $Y$ are discrete random variables with joint PMF $p(x_i,y_j)$, then

\[
F(x,y)=\sum_{x_i\le x}\sum_{y_j\le y}p(x_i,y_j)
\]

Interpretation:

To compute the joint CDF:

\begin{itemize}
\item Add probabilities from the joint PMF table
\item Include rows where $x_i\le x$
\item Include columns where $y_j\le y$
\end{itemize}

\section{Properties of the Joint CDF}

The joint CDF $F(x,y)$ satisfies the following properties:

\subsection*{1. Monotonicity}

$F(x,y)$ is non-decreasing.

If $x$ or $y$ increases, $F(x,y)$ remains constant or increases.

\subsection*{2. Lower Boundary Condition}

\[
F(x,y)=0
\]

at the lower-left corner of the joint range.

If the lower bound is $(-\infty,-\infty)$, then

\[
\lim_{(x,y)\to(-\infty,-\infty)}F(x,y)=0
\]

\subsection*{3. Upper Boundary Condition}

\[
F(x,y)=1
\]

at the upper-right corner of the joint range.

If the upper bound is $(\infty,\infty)$, then

\[
\lim_{(x,y)\to(\infty,\infty)}F(x,y)=1
\]

\section{Summary of Lecture}

This lecture introduced the concept of joint random variables and joint distributions.

Main ideas covered:

\begin{itemize}
\item Motivation for studying joint random variables
\item Joint PMF for discrete variables
\item Joint PMF table representation
\item Example using two dice
\item Joint PDF for continuous variables
\item Example computation with a joint PDF
\item Definition of the joint CDF
\item Continuous and discrete forms of joint CDF
\item Properties of the joint CDF
\end{itemize}

These concepts form the foundation for analyzing relationships between multiple random variables in probabilistic systems.

\end{document}